% Part: lambda-calculus
% Chapter: church-rosser
% Section: beta-reduction

\documentclass[../../../include/open-logic-section]{subfiles}

\begin{document}

\olfileid{lam}{cr}{b}

\olsection{$\beta$-લઘુકરણ}

\begin{lem}\ollabel{lem:one-par}
  $M \bredone M'$ હોય તો $M \bredpar M'$.
\end{lem}
\begin{proof}
  $M \bredone M'$ની રચના પર આગમનથી. મૂળ સંકોચનનો
  કિસ્સામાં $M$ એ $(\lambd[x][N])Q$ અને
  $M'$ એ $\Subst{N}{Q}{x}$ છે, કોઈક
  $x$, $N$ અને~$Q$ માટે. \olref[pb]{thm:refl}
  મુજબ $N \bredpar N$ અને $Q \bredpar Q$;
  તેથી \olref[pb]{defn:bredpar}\olref[pb]{defn:bredpar4}
  મુજબ $(\lambd[x][N])Q \bredpar
  \Subst{N}{Q}{x}$. જો સંકોચન અમૂર્તીકરણની અંદર
  થયું હોય તો આગમનની ધારણા અને
  \olref[pb]{defn:bredpar2} વાપરીએ. અનુપ્રયોગની
  ડાબી અથવા જમણી બાજુ થયું હોય તો બદલાતી બાજુ
  માટે આગમનની ધારણા, બીજી માટે
  \olref[pb]{thm:refl} અને પછી
  \olref[pb]{defn:bredpar3} વાપરીએ.
  \sourcecorrection{OLLAM-036}{સ્થિર અંગ્રેજી મૂળની
  આ સાબિતી ફક્ત આખા પદના મૂળ સંકોચનનો કિસ્સો
  લે છે, જ્યારે $\bredone$ સૌથી નાનો
  રચના-સુસંગત સંબંધ છે અને સંકોચન અંદર પણ થઈ શકે.
  અહીં અમૂર્તીકરણ અને અનુપ્રયોગ હેઠળના બાકીના
  આગમનાત્મક કિસ્સા ઉમેર્યા છે.}
\end{proof}

\begin{lem}\ollabel{lem:par-red}
  $M \bredpar M'$ હોય તો $M \bred M'$.
\end{lem}

\begin{proof}
  $M \bredpar M'$ની !!{derivation} પર આગમનથી.
  \begin{enumerate}
  \item છેલ્લો નિયમ \olref[pb]{defn:bredpar1} હોય:
    $M$ અને $M'$ બંને ચલ~$x$ છે, તેથી $x \bred x$.
  \item છેલ્લો નિયમ \olref[pb]{defn:bredpar2} હોય:
    $M$ એ $\lambd[x][N]$ અને $M'$ એ
    $\lambd[x][N']$, કોઈક $x$, $N$, $N'$ માટે,
    જ્યાં $N \bredpar N'$. આગમનની ધારણાથી
    $N \bred N'$. $N \bred N'$ની જ $\bredone$-સંકોચન શ્રેણી
    અમૂર્તીકરણની અંદર
    લાગુ કરતાં $\lambd[x][N] \bred
    \lambd[x][N']$ મળે છે.
  \item છેલ્લો નિયમ \olref[pb]{defn:bredpar3} હોય:
    $M$ એ $PQ$ અને $M'$ એ $P'Q'$ છે,
    કોઈક $P$, $Q$, $P'$, $Q'$ માટે, જ્યાં
    $P \bredpar P'$ અને $Q \bredpar Q'$.
    આગમનની ધારણાથી $P \bred P'$ અને
    $Q \bred Q'$. પહેલાં $P \bred P'$,
    પછી $Q \bred Q'$ મુજબ લઘુકરણ કરીએ તો
    $PQ \bred P'Q'$ મળે છે.
  \item છેલ્લો નિયમ \olref[pb]{defn:bredpar4} હોય:
    $M$ એ $(\lambd[x][N])Q$ અને
    $M'$ એ $\Subst{N'}{Q'}{x}$, કોઈક
    $x$, $N$, $N'$, $Q$, $Q'$ માટે; અહીં
    $N \bredpar N'$ તથા $Q \bredpar Q'$.
    આગમનની ધારણાથી $Q \bred Q'$ અને
    $N \bred N'$. પહેલાં $N \bred N'$, પછી
    $Q \bred Q'$ મુજબ લઘુકરણ કરીને છેલ્લે
    $(\lambd[x][N'])Q'$નું
    $\Subst{N'}{Q'}{x}$માં સંકોચન કરીએ તો
    $(\lambd[x][N])Q \bred
    \Subst{N'}{Q'}{x}$ મળે છે.
    \sourcecorrection{OLLAM-037}{સ્થિર અંગ્રેજી મૂળ
    આ છેલ્લી કિસ્સાની પદ-યાદીમાં $N'$ની જગ્યાએ
    ફરી~$M'$ લખે છે. અહીં નિયમની પૂર્વધારણા પ્રમાણે
    $N'$ પુનઃસ્થાપિત કર્યું છે.}
  \end{enumerate}
\end{proof}

\begin{lem}\ollabel{lem:str}
  $\bred$ એ $\bredpar$ને સમાવતો સૌથી નાનો
  પરંપરિત સંબંધ છે.
\end{lem}

\begin{proof}
  $\xred$ એ $\bredpar$ને સમાવતો સૌથી નાનો
  પરંપરિત સંબંધ માનો.

  $\bred \subseteq \xred$: માનો કે $M \bred M'$.
  શૂન્ય પગલાંના કિસ્સામાં $M=M'$ છે, અને
  \olref[pb]{thm:refl}થી $M \bredpar M$, તેથી
  $M \xred M'$. બાકી $M \ident M_1
  \bredone \dots \bredone M_k \ident M'$.
  \olref{lem:one-par}થી $M \ident M_1
  \bredpar \dots \bredpar M_k \ident M'$.
  $\xred$ એ $\bredpar$ને સમાવતો પરંપરિત
  સંબંધ હોવાથી $M \xred M'$.
  \sourcecorrection{OLLAM-038}{સ્થિર અંગ્રેજી મૂળ
  $\bred$ના શૂન્ય-પગલાના સ્વવાચક કિસ્સાને
  એક-અથવા-વધુ પગલાંની શ્રેણીથી આવરી લેવાનો
  પ્રયાસ કરે છે. અહીં પહેલાં $\bredpar$ની
  સ્વવાચકતાથી તે કિસ્સો અલગ સાબિત કર્યો છે.}

  $\xred \subseteq \bred$: માનો કે $M \xred M'$,
  એટલે $M \ident M_1 \bredpar \dots
  \bredpar M_k \ident M'$. \olref{lem:par-red}થી
  $M \ident M_1 \bred \dots \bred
  M_k \ident M'$. $\bred$ પરંપરિત હોવાથી
  $M \bred M'$.
\end{proof}

\begin{thm}\ollabel{thm:cr}
  $\bred$ ચર્ચ--રોસર ગુણધર્મ ધરાવે છે.
\end{thm}

\begin{proof}
  \olref[dap]{thm:str}, \olref[pb]{thm:cr} અને
  \olref{lem:str}માંથી તરત મળે છે. જોકે
  \olref[pb]{thm:cr}ની અગાઉ નોંધેલી સાબિતી-ખામી
  હજી ખુલ્લી છે.
\end{proof}

\end{document}
